\chapter{RESULTS}
\label{chapter4}
This chapter presents the results obtained from the study. The results are organized in a structured manner to clearly demonstrate the findings derived from experiments, simulations, or analytical procedures. The purpose of this chapter is to provide factual and objective outcomes without interpretation, which will be discussed in the next chapter.

This chapter is divided into different sections based on the objectives of the study. Each section presents results corresponding to specific aspects of the research.

% -------------------------------------------------
\section{Overview of Results}

This section provides a general overview of the results. The findings are presented using tables, figures, and charts where appropriate. The results should be clear, concise, and directly related to the research objectives.

Figure \ref{fig:sample_result_plot} shows an example of graphical representation of experimental data.

\begin{figure}[!tbh]
    \centering
    \rule{0.7\textwidth}{5cm}
    \vspace{0.5em}
    \caption{Example plot showing variation of output parameter with respect to input condition.}
    \label{fig:sample_result_plot}
\end{figure}

% -------------------------------------------------
\section{Topic 1 Results}

This section presents the results related to the first objective or topic of the study.

\subsection{Case 1}

Describe the results obtained for Case 1. Provide detailed observations supported by data.

\begin{figure}[!tbh]
    \centering
    \rule{0.65\textwidth}{4.5cm}
    \vspace{0.5em}
    \caption{Result visualization for Topic 1 - Case 1.}
    \label{fig:topic1_case1}
\end{figure}

\subsection{Case 2}

Describe the results obtained for Case 2. Highlight any variations or differences observed compared to Case 1.

\begin{table}[!tbh]
\centering
\caption{Comparison of results for different cases under Topic 1}
\label{tab:topic1_cases}
\begin{tabular}{|c|c|c|}
\hline
\textbf{Case} & \textbf{Parameter Value} & \textbf{Output Result} \\
\hline
Case 1 & Value 1 & Result 1 \\
Case 2 & Value 2 & Result 2 \\
\hline
\end{tabular}
\vspace{0.3em}
\end{table}

% -------------------------------------------------
\section{Topic 2 Results}

This section presents the results related to the second objective or topic.

\subsection{Case 1}

Provide details of the observed results. Include graphs, images, or numerical data where necessary.

\begin{figure}[!tbh]
    \centering
    \rule{0.65\textwidth}{4.5cm}
    \vspace{0.5em}
    \caption{Result visualization for Topic 2 - Case 1.}
    \label{fig:topic2_case1}
\end{figure}

\subsection{Case 2}

Present the results for Case 2 and compare it with Case 1 if applicable.

% -------------------------------------------------
\section{Topic 3 Results}

This section presents results related to the third objective or research component.

\subsection{Case 1}

Provide the data obtained and explain the observations.

\subsection{Case 2}

Present comparative or additional results.

% -------------------------------------------------
\section{Summary of Results}

This section provides a summary of key findings from all sections of this chapter. The summary should highlight important trends, variations, or patterns observed in the results.

For example:
\begin{itemize}
    \item Key trend observed in the data
    \item Comparison between different cases
    \item Notable patterns or anomalies
\end{itemize}

The results presented in this chapter will be further analyzed and interpreted in the next chapter.